\chapter{Chapter09 — KKL Bribery}
%%% generated by the blueprint pipeline from FABL.Chapter09.KKLBribery
\section{Overview}

Declarations of the Lean module \texttt{FABL.Chapter09.KKLBribery}.

\section{Declarations}

\begin{definition}[Bribed input]
\lean{FABL.bribeInput}
\statementsource{Analysis-of-Boolean-Functions-by-Ryan-ODonnell}{pdf-ffbe9c7b3306-p265-b007}
\label{FABL.bribeInput}
\leanok
\uses{FABL.Sign, FABL.SignCube}
For a coordinate set $J\subseteq[n]$, a sign $b\in\{-1,1\}$, and an input
$x\in\{-1,1\}^n$, define the modified input by setting every coordinate in $J$ equal to
$b$ and leaving every coordinate outside $J$ equal to its value in $x$.
\end{definition}

\begin{definition}[Bribed Boolean function]
\lean{FABL.bribeFunction}
\statementsource{Analysis-of-Boolean-Functions-by-Ryan-ODonnell}{pdf-ffbe9c7b3306-p265-b007}
\label{FABL.bribeFunction}
\leanok
\uses{FABL.BooleanFunction, FABL.Sign, FABL.bribeInput}
For a Boolean function $f:\{-1,1\}^n\to\{-1,1\}$, a coordinate set $J\subseteq[n]$, and
a sign $b\in\{-1,1\}$, define the modified function by evaluating $f$ on the input
obtained by setting the coordinates in $J$ to $b$ and leaving all other coordinates
unchanged.
\end{definition}

\begin{theorem}[Bribing no coordinates]
\lean{FABL.bribeFunction_empty}
\label{FABL.bribeFunction_empty}
\leanok
\uses{FABL.BooleanFunction, FABL.Sign, FABL.bribeFunction, FABL.bribeInput}
\difficulty{3}
Let $n$ be a natural number, let $f$ be a Boolean function on the $n$-dimensional sign
cube, and let $b\in\{-1,1\}$. Replacing the coordinates in the empty set by $b$ leaves
$f$ unchanged.
\end{theorem}

\begin{theorem}[Bribing an inserted coordinate]
\lean{FABL.bribeInput_insert}
\label{FABL.bribeInput_insert}
\leanok
\uses{FABL.Sign, FABL.SignCube, FABL.bribeInput, FABL.setCoordinate, FABL.setCoordinate_apply_of_ne}
\difficulty{4}
Let $n$ be a nonnegative integer, let $J$ be a finite set of indices in
$\{0,\ldots,n-1\}$, let $i$ be such an index, and let $b\in\{-1,1\}$. For every sign
vector $x\in\{-1,1\}^n$, replacing the coordinates in $J\cup\{i\}$ by $b$ gives the same
vector as first setting coordinate $i$ of $x$ to $b$ and then replacing the coordinates
in $J$ by $b$.
\end{theorem}

\begin{theorem}[Bribing an enlarged coordinate set]
\lean{FABL.bribeFunction_insert}
\label{FABL.bribeFunction_insert}
\leanok
\uses{FABL.BooleanFunction, FABL.Sign, FABL.bribeFunction, FABL.bribeInput_insert, FABL.setCoordinate}
\difficulty{3}
Let $n$ be a natural number, let $f$ be a Boolean function on the $n$-dimensional sign
cube, let $J$ be a set of coordinates, let $i$ be a coordinate, and let $b\in\{-1,1\}$.
For every input $x$, bribing the coordinates in $J\cup\{i\}$ to $b$ gives the same
output as first setting coordinate $i$ of $x$ to $b$ and then bribing the coordinates in
$J$ to $b$.
\end{theorem}

\begin{theorem}[Monotonicity under coordinate fixing]
\lean{FABL.monotone_bribeFunction}
\label{FABL.monotone_bribeFunction}
\leanok
\uses{FABL.BooleanFunction, FABL.Sign, FABL.bribeFunction, FABL.bribeInput}
\difficulty{4}
Let $f$ be a monotone Boolean function on the $n$-dimensional sign cube, let $J$ be a
set of coordinates, and fix a sign $b\in\{-1,1\}$. The function obtained by setting
every coordinate in $J$ to $b$ and leaving all other coordinates unchanged before
evaluating $f$ is monotone.
\end{theorem}

\begin{theorem}[Influence of a bribed coordinate]
\lean{FABL.booleanInfluence_bribeFunction_eq_zero_of_mem}
\label{FABL.booleanInfluence_bribeFunction_eq_zero_of_mem}
\leanok
\uses{FABL.BooleanFunction, FABL.IsPivotal, FABL.Sign, FABL.booleanInfluence, FABL.bribeFunction, FABL.bribeInput, FABL.flipCoordinate, FABL.uniformProbability}
\difficulty{5}
Let $n$ be a nonnegative integer, let $f$ be a Boolean function on the $n$-dimensional
sign cube, let $J$ be a set of coordinates, and choose a sign $b$. If $i$ belongs to
$J$, then the influence of the function obtained by setting every coordinate in $J$ to
$b$ before applying $f$ is zero in coordinate $i$.
\end{theorem}

\begin{theorem}[Mean of a sign-valued function]
\lean{FABL.mean_toReal_mem_Icc}
\label{FABL.mean_toReal_mem_Icc}
\leanok
\uses{FABL.BooleanFunction, FABL.BooleanFunction.toReal, FABL.SignCube, FABL.mean}
\difficulty{4}
Let $n$ be a nonnegative integer and let $f$ map the sign cube $\{−1,1\}^n$ to the two
signs. The mean of the real-valued function obtained by viewing each output sign as the
real number $-1$ or $1$ lies in $[-1,1]$.
\end{theorem}

\begin{theorem}[Mean under coordinate averaging]
\lean{FABL.mean_coordinateExpectation}
\label{FABL.mean_coordinateExpectation}
\leanok
\uses{FABL.SignCube, FABL.coordinateExpectation, FABL.coordinateLaplacian, FABL.fourierCoeff_coordinateLaplacian, FABL.mean, FABL.mean_eq_fourierCoeff_empty}
\difficulty{4}
Let $f$ be a real-valued function on the sign cube of dimension $n$, and choose any
coordinate $i$. Replace the value of $f$ at each input by the average of its values when
coordinate $i$ is set to $+1$ and to $-1$. This coordinate average has the same mean as
$f$.
\end{theorem}

\begin{theorem}[Mean after fixing a coordinate]
\lean{FABL.mean_fixCoordinate_one}
\label{FABL.mean_fixCoordinate_one}
\leanok
\uses{FABL.SignCube, FABL.coordinateExpectation, FABL.coordinateExpectation_apply, FABL.discreteDerivative, FABL.discreteDerivative_apply, FABL.fourierCoeff, FABL.mean, FABL.mean_coordinateExpectation, FABL.mean_discreteDerivative_eq_fourierCoeff_singleton, FABL.setCoordinate}
\difficulty{4}
Let $n$ be a nonnegative integer, let $f:\{−1,1\}^n\to\mathbb R$, and let $i$ be a
coordinate. The uniform mean of $f$ with coordinate $i$ fixed to $1$ equals the uniform
mean of $f$ plus its Fourier coefficient at the singleton set $\{i\}$.
\end{theorem}

\begin{theorem}[Mean change from adding a coordinate]
\lean{FABL.mean_bribeFunction_insert_one}
\label{FABL.mean_bribeFunction_insert_one}
\leanok
\uses{FABL.BooleanFunction, FABL.BooleanFunction.toReal, FABL.booleanInfluence, FABL.booleanInfluence_eq_influence_toReal, FABL.bribeFunction, FABL.bribeFunction_insert, FABL.influence_eq_fourierCoeff_singleton_of_monotone, FABL.mean, FABL.mean_fixCoordinate_one, FABL.monotone_bribeFunction, FABL.setCoordinate}
\difficulty{3}
Let $f$ be a monotone Boolean function on the $n$-dimensional sign cube, let $J$ be a
set of coordinates, and let $i$ be a coordinate. The mean after replacing all
coordinates in $J\cup\{i\}$ by $1$ equals the mean after replacing the coordinates in
$J$ by $1$, plus the influence of coordinate $i$ on the resulting function.
\end{theorem}

\begin{definition}[Bribery increment]
\lean{FABL.briberyIncrement}
\statementsource{Analysis-of-Boolean-Functions-by-Ryan-ODonnell}{pdf-ffbe9c7b3306-p266-b001}
\label{FABL.briberyIncrement}
\leanok
For each natural number $n$, define the bribery increment by
$199\log(n)/(1{,}000{,}000n)$, where $\log$ is the natural logarithm.
\end{definition}

\begin{theorem}[Variance within a mean window]
\lean{FABL.variance_ge_of_mean_mem_bribery_window}
\proofsource{Analysis-of-Boolean-Functions-by-Ryan-ODonnell}{pdf-ffbe9c7b3306-p266-b001}
\label{FABL.variance_ge_of_mean_mem_bribery_window}
\leanok
\uses{FABL.BooleanFunction, FABL.BooleanFunction.toReal, FABL.mean, FABL.variance, FABL.variance_eq_four_mul_probabilities}
\difficulty{3}
Let $n$ be a nonnegative integer and let $f$ be a function from the $n$-dimensional sign
cube to $\{ -1,1\}$. If its uniform mean satisfies $-99/100\leq \mathbb{E}f\leq 99/100$,
then its variance under the uniform distribution is at least $199/10000$.
\end{theorem}

\begin{theorem}[A fresh bribery step]
\lean{FABL.exists_fresh_bribery_step}
\proofsource{Analysis-of-Boolean-Functions-by-Ryan-ODonnell}{pdf-ffbe9c7b3306-p265-b007}
\label{FABL.exists_fresh_bribery_step}
\leanok
\uses{FABL.BooleanFunction, FABL.BooleanFunction.toReal, FABL.booleanInfluence, FABL.booleanInfluence_bribeFunction_eq_zero_of_mem, FABL.bribeFunction, FABL.briberyIncrement, FABL.exists_booleanInfluence_eq_maximumInfluence, FABL.kkl, FABL.maximumInfluence, FABL.mean, FABL.mean_bribeFunction_insert_one, FABL.variance, FABL.variance_ge_of_mean_mem_bribery_window}
\difficulty{5}
Let $n\ge 2$, let $f$ be a monotone Boolean function on the $n$-dimensional sign cube,
and let $J$ be a set of coordinates. Suppose the mean of the function obtained by
setting every coordinate in $J$ to $+1$ lies in $[-99/100,99/100)$. Then some coordinate
$i\notin J$ can be added to $J$ so that this mean increases by at least
$199\log(n)/(1000000n)$.
\end{theorem}

\begin{definition}[Maximum influence coordinate]
\lean{FABL.maximumInfluenceCoordinate}
\statementsource{Analysis-of-Boolean-Functions-by-Ryan-ODonnell}{pdf-ffbe9c7b3306-p265-b003}
\label{FABL.maximumInfluenceCoordinate}
\leanok
\uses{FABL.BooleanFunction, FABL.exists_booleanInfluence_eq_maximumInfluence}
For a positive integer $n$ and a Boolean function $f:\{-1,1\}^n\to\{-1,1\}$, let
$i_f\in\{1,\ldots,n\}$ be the coordinate selected to have maximum influence among the
coordinates of $f$.
\end{definition}

\begin{theorem}[Coordinate attaining maximum influence]
\lean{FABL.booleanInfluence_maximumInfluenceCoordinate}
\label{FABL.booleanInfluence_maximumInfluenceCoordinate}
\leanok
\uses{FABL.BooleanFunction, FABL.booleanInfluence, FABL.exists_booleanInfluence_eq_maximumInfluence, FABL.maximumInfluence, FABL.maximumInfluenceCoordinate}
\difficulty{1}
Let $n$ be a positive integer, let $f:\{-1,1\}^n\to\{\pm1\}$, and let $i_*$ be the
coordinate selected as a maximum-influence coordinate of $f$. Its influence equals the
maximum influence of $f$.
\end{theorem}

\begin{theorem}[Bribery step using a maximum-influence coordinate]
\lean{FABL.maximumInfluenceCoordinate_bribery_step}
\proofsource{Analysis-of-Boolean-Functions-by-Ryan-ODonnell}{pdf-ffbe9c7b3306-p266-b001}
\label{FABL.maximumInfluenceCoordinate_bribery_step}
\leanok
\uses{FABL.BooleanFunction, FABL.BooleanFunction.toReal, FABL.booleanInfluence, FABL.booleanInfluence_bribeFunction_eq_zero_of_mem, FABL.booleanInfluence_le_maximumInfluence, FABL.booleanInfluence_maximumInfluenceCoordinate, FABL.bribeFunction, FABL.briberyIncrement, FABL.exists_fresh_bribery_step, FABL.maximumInfluence, FABL.maximumInfluenceCoordinate, FABL.mean, FABL.mean_bribeFunction_insert_one}
\difficulty{5}
Let $n\ge 2$, let $f:\{-1,1\}^n\to\{-1,1\}$ be monotone, and let
$J\subseteq\{0,\ldots,n-1\}$. Suppose the mean of the function obtained by setting all
coordinates in $J$ to $1$ lies in $[-99/100,99/100)$. Choose a coordinate $i$ of maximum
influence for that modified function. Then $i\notin J$, and setting the coordinates in
$J\cup\{i\}$ to $1$ increases the mean by at least $199\log(n)/(1{,}000{,}000n)$.
\end{theorem}

\begin{definition}[Iterated bribery sets]
\lean{FABL.briberySets}
\statementsource{Analysis-of-Boolean-Functions-by-Ryan-ODonnell}{pdf-ffbe9c7b3306-p265-b007}
\label{FABL.briberySets}
\leanok
\uses{FABL.BooleanFunction, FABL.BooleanFunction.toReal, FABL.bribeFunction, FABL.maximumInfluenceCoordinate, FABL.mean}
Let $n>0$ and let $f:\{−1,1\}^n\to\{−1,1\}$. Define a sequence of coordinate sets by
$J_0=\varnothing$. Given $J_t$, form the function obtained by setting every coordinate
in $J_t$ to $1$ before evaluating $f$. If its uniform mean on $\{−1,1\}^n$ is less than
$99/100$, let $J_{t+1}$ be $J_t$ together with a coordinate of maximum influence for
that modified function; otherwise let $J_{t+1}=J_t$.
\end{definition}

\begin{theorem}[Initial bribery set]
\lean{FABL.briberySets_zero}
\label{FABL.briberySets_zero}
\leanok
\uses{FABL.BooleanFunction, FABL.briberySets}
\difficulty{1}
Let $n>0$ and let $f$ be a Boolean function on the $n$-dimensional sign cube. The
bribery set at stage $0$ is empty.
\end{theorem}

\begin{theorem}[Successive bribery sets]
\lean{FABL.briberySets_succ}
\proofsource{Analysis-of-Boolean-Functions-by-Ryan-ODonnell}{pdf-ffbe9c7b3306-p265-b007}
\label{FABL.briberySets_succ}
\leanok
\uses{FABL.BooleanFunction, FABL.BooleanFunction.toReal, FABL.bribeFunction, FABL.briberySets, FABL.maximumInfluenceCoordinate, FABL.mean}
\difficulty{1}
Let $n>0$ and let $f$ be a Boolean function on the $n$-dimensional sign cube. For every
natural number $t$, the bribery set at stage $t+1$ is obtained from the stage-$t$ set by
adjoining a coordinate of maximum influence for the function bribed on that set to
output $1$ if the mean of that bribed function is less than $99/100$; otherwise, the set
is unchanged.
\end{theorem}

\begin{theorem}[Upper bound on the bribery set size]
\lean{FABL.card_briberySets_le}
\label{FABL.card_briberySets_le}
\leanok
\uses{FABL.BooleanFunction, FABL.briberySets, FABL.briberySets_succ}
\difficulty{4}
Let $n>0$, let $f:\{-1,1\}^n\to\{-1,1\}$, and let $t\in\mathbb N$. The bribery
set formed after $t$ steps, with target $99/100$, has cardinality at most $t$.
\end{theorem}

\begin{theorem}[Nonnegativity of the bribery increment]
\lean{FABL.briberyIncrement_nonneg}
\label{FABL.briberyIncrement_nonneg}
\leanok
\uses{FABL.briberyIncrement}
\difficulty{4}
For every nonnegative integer $n$, the quantity $\frac{199\log n}{1000000n}$, with the
logarithm and arithmetic interpreted over the real numbers, is nonnegative.
\end{theorem}

\begin{theorem}[Mean gain from bribery sets]
\lean{FABL.mean_briberySets_ge_initial_add_card_mul}
\proofsource{Analysis-of-Boolean-Functions-by-Ryan-ODonnell}{pdf-ffbe9c7b3306-p266-b001}
\label{FABL.mean_briberySets_ge_initial_add_card_mul}
\leanok
\uses{FABL.BooleanFunction, FABL.BooleanFunction.toReal, FABL.bribeFunction, FABL.briberyIncrement, FABL.briberyIncrement_nonneg, FABL.briberySets, FABL.briberySets_succ, FABL.maximumInfluenceCoordinate, FABL.maximumInfluenceCoordinate_bribery_step, FABL.mean}
\difficulty{5}
Let $f$ be a coordinatewise monotone Boolean function on the sign cube of dimension $n$,
where $n\ge 2$, and suppose its mean is at least $-99/100$. For every natural number
$t$, forcing the coordinates in the bribery set obtained after $t$ steps to $+1$ gives a
function with mean at least the original mean plus the size of that set times
$199\log(n)/(1{,}000{,}000n)$.
\end{theorem}

\begin{theorem}[Bribery set size under a low mean]
\lean{FABL.card_briberySets_eq_of_mean_lt_target}
\proofsource{Analysis-of-Boolean-Functions-by-Ryan-ODonnell}{pdf-ffbe9c7b3306-p266-b001}
\label{FABL.card_briberySets_eq_of_mean_lt_target}
\leanok
\uses{FABL.BooleanFunction, FABL.BooleanFunction.toReal, FABL.bribeFunction, FABL.briberyIncrement_nonneg, FABL.briberySets, FABL.briberySets_succ, FABL.maximumInfluenceCoordinate, FABL.maximumInfluenceCoordinate_bribery_step, FABL.mean, FABL.mean_briberySets_ge_initial_add_card_mul}
\difficulty{5}
Let $n\ge 2$ and let $f:\{-1,1\}^n\to\{-1,1\}$ be monotone. Suppose its mean is at
least $-99/100$, and let $t\in\mathbb N$ satisfy $t\le n$. If the function obtained
after setting to $1$ all coordinates in the bribery set after $t$ steps has mean less
than $99/100$, then that bribery set has cardinality exactly $t$.
\end{theorem}

\begin{theorem}[Mean bribe function univ one]
\lean{FABL.mean_bribeFunction_univ_one}
\statementsource{Analysis-of-Boolean-Functions-by-Ryan-ODonnell}{pdf-ffbe9c7b3306-p266-b001}
\label{FABL.mean_bribeFunction_univ_one}
\leanok
\uses{FABL.BooleanFunction, FABL.BooleanFunction.toReal, FABL.Sign, FABL.bribeFunction, FABL.bribeInput, FABL.mean, FABL.neg_one_le_sign, FABL.sign_le_one}
\difficulty{4}
% dataset-pipeline: stub -- prose pending build_informal_statements.py
Statement pending regeneration.
\end{theorem}

\begin{theorem}[Final mean after iterative bribery]
\lean{FABL.mean_briberySets_final_ge_target}
\proofsource{Analysis-of-Boolean-Functions-by-Ryan-ODonnell}{pdf-ffbe9c7b3306-p266-b001}
\label{FABL.mean_briberySets_final_ge_target}
\leanok
\uses{FABL.BooleanFunction, FABL.BooleanFunction.toReal, FABL.bribeFunction, FABL.briberySets, FABL.card_briberySets_eq_of_mean_lt_target, FABL.mean, FABL.mean_bribeFunction_univ_one}
\difficulty{3}
Let $f$ be a coordinatewise monotone Boolean function on the sign cube of dimension $n$,
where $n\ge 2$, and suppose its mean is at least $-99/100$. After the bribery procedure
has run for $n$ steps, forcing the coordinates in its resulting set to $+1$ gives a
function whose mean is at least $99/100$.
\end{theorem}

\begin{theorem}[Proposition9 27 positive]
\lean{FABL.proposition9_27_positive}
\statementsource{Analysis-of-Boolean-Functions-by-Ryan-ODonnell}{
  pdf-ffbe9c7b3306-p265-b006,
  pdf-ffbe9c7b3306-p265-b007,
  pdf-ffbe9c7b3306-p266-b001,
  pdf-ffbe9c7b3306-p097-b006
}
\label{FABL.proposition9_27_positive}
\leanok
\uses{FABL.BooleanFunction, FABL.BooleanFunction.toReal, FABL.bribeFunction, FABL.briberyIncrement, FABL.briberySets, FABL.mean, FABL.mean_briberySets_final_ge_target, FABL.mean_briberySets_ge_initial_add_card_mul, FABL.mean_toReal_mem_Icc}
\difficulty{4}
% dataset-pipeline: stub -- prose pending build_informal_statements.py
Statement pending regeneration.
\end{theorem}

\begin{definition}[Dual Boolean function]
\lean{FABL.dualBooleanFunction}
\statementsource{Analysis-of-Boolean-Functions-by-Ryan-ODonnell}{pdf-ffbe9c7b3306-p094-b005}
\label{FABL.dualBooleanFunction}
\leanok
\uses{FABL.BooleanFunction}
For a Boolean function $f:\{−1,1\}^n\to\{−1,1\}$, define its dual by $f^*(x)=-f(-x)$,
where $-x$ negates every coordinate of $x$.
\end{definition}

\begin{definition}[Coordinatewise sign negation]
\lean{FABL.briberyNegateSignCubeEquiv}
\label{FABL.briberyNegateSignCubeEquiv}
\leanok
\uses{FABL.SignCube}
For each natural number $n$, the map on the $n$-dimensional sign cube that sends each
sign vector $x$ to the vector whose $i$th coordinate is $-x_i$ is a bijection.
\end{definition}

\begin{theorem}[Monotonicity of the dual function]
\lean{FABL.monotone_dualBooleanFunction}
\label{FABL.monotone_dualBooleanFunction}
\leanok
\uses{FABL.BooleanFunction, FABL.dualBooleanFunction, FABL.neg_one_le_sign}
\difficulty{5}
Let $f$ be a monotone Boolean function on the $n$-dimensional sign cube. The function
$x\mapsto -f(-x)$ is monotone on that cube.
\end{theorem}

\begin{theorem}[Mean of the dual Boolean function]
\lean{FABL.mean_dualBooleanFunction}
\label{FABL.mean_dualBooleanFunction}
\leanok
\uses{FABL.BooleanFunction, FABL.BooleanFunction.toReal, FABL.SignCube, FABL.briberyNegateSignCubeEquiv, FABL.dualBooleanFunction, FABL.mean, FABL.signValue}
\difficulty{3}
Let $f$ be a Boolean function on the sign cube of dimension $n$. Define its dual by
negating the output of $f$ after negating every input coordinate. The mean of this dual
function is the negative of the mean of $f$.
\end{theorem}

\begin{theorem}[Mean under dual bribery]
\lean{FABL.mean_bribeFunction_dual_one}
\label{FABL.mean_bribeFunction_dual_one}
\leanok
\uses{FABL.BooleanFunction, FABL.BooleanFunction.toReal, FABL.SignCube, FABL.bribeFunction, FABL.bribeInput, FABL.briberyNegateSignCubeEquiv, FABL.dualBooleanFunction, FABL.mean, FABL.signValue}
\difficulty{4}
Let $f$ be a Boolean function on the $n$-dimensional sign cube and let $J$ be a set of
coordinates. The mean of the function obtained by first replacing the coordinates in $J$
by $1$ and then applying the dual of $f$ equals the negative of the mean obtained by
replacing the coordinates in $J$ by $-1$ and applying $f$.
\end{theorem}

\begin{theorem}[A small coalition forcing a negative bias]
\lean{FABL.proposition9_27_negative}
\proofsource{Analysis-of-Boolean-Functions-by-Ryan-ODonnell}{
  pdf-ffbe9c7b3306-p265-b006,
  pdf-ffbe9c7b3306-p097-b006
}
\label{FABL.proposition9_27_negative}
\leanok
\uses{FABL.BooleanFunction, FABL.BooleanFunction.toReal, FABL.bribeFunction, FABL.dualBooleanFunction, FABL.mean, FABL.mean_bribeFunction_dual_one, FABL.mean_dualBooleanFunction, FABL.monotone_dualBooleanFunction, FABL.proposition9_27_positive}
\difficulty{3}
Let $n\ge 2$, and let $f:\{-1,1\}^n\to\{-1,1\}$ be monotone, with the signs ordered by
$-1<1$. If the uniform mean of $f$ is at most $99/100$, there is a set $J$ of
coordinates of size at most $10000n/\log n$ such that forcing every coordinate in $J$ to
$-1$ gives a function whose uniform mean is at most $-99/100$.
\end{theorem}

\begin{theorem}[Small coalitions can bias a monotone Boolean function]
\lean{FABL.proposition9_27}
\proofsource{Analysis-of-Boolean-Functions-by-Ryan-ODonnell}{
  pdf-ffbe9c7b3306-p265-b006,
  pdf-ffbe9c7b3306-p265-b007,
  pdf-ffbe9c7b3306-p266-b001
}
\label{FABL.proposition9_27}
\leanok
\uses{FABL.BooleanFunction, FABL.BooleanFunction.toReal, FABL.bribeFunction, FABL.mean, FABL.proposition9_27_negative, FABL.proposition9_27_positive}
\difficulty{1}
Let $n\ge 2$, and let $f:\{-1,1\}^n\to\{-1,1\}$ be monotone, with the signs ordered by
$-1<1$. If the uniform mean of $f$ is at least $-99/100$, there is a set $J$ of
coordinates such that forcing every coordinate in $J$ to $1$ gives a function with mean
at least $99/100$, and $|J|\le 10000n/\log n$. If the mean of $f$ is at most $99/100$,
there is a set $J$ such that forcing every coordinate in $J$ to $-1$ gives a function
with mean at most $-99/100$, and $|J|\le 10000n/\log n$.
\end{theorem}

\begin{definition}[Coalition can attain an output]
\lean{FABL.CoalitionCanAttain}
\statementsource{Analysis-of-Boolean-Functions-by-Ryan-ODonnell}{
  pdf-ffbe9c7b3306-p016-b005,
  pdf-ffbe9c7b3306-p274-b007,
  pdf-ffbe9c7b3306-p275-b001
}
\label{FABL.CoalitionCanAttain}
\leanok
\uses{FABL.BooleanFunction, FABL.Sign, FABL.SignCube}
Let $f:\{−1,1\}^n\to\{−1,1\}$, let $J\subseteq\{1,\ldots,n\}$ be a set of coordinates,
let $b\in\{−1,1\}$, and let $x\in\{−1,1\}^n$. The coalition of coordinates in $J$ can
attain output $b$ from $x$ if there is an input $y$ agreeing with $x$ outside $J$ such
that $f(y)=b$.
\end{definition}

\begin{definition}[Coalition success probability]
\lean{FABL.coalitionSuccessProbability}
\label{FABL.coalitionSuccessProbability}
\leanok
\uses{FABL.BooleanFunction, FABL.CoalitionCanAttain, FABL.Sign, FABL.uniformProbability}
For a Boolean function $f$ on the $n$-dimensional sign cube, a set $J$ of coordinates,
and a sign $b\in\{-1,1\}$, the coalition success probability is the uniform probability
that the input satisfies the condition that the coalition on $J$ can attain $b$ under
$f$.
\end{definition}

\begin{definition}[Coalitional influence toward an output]
\lean{FABL.coalitionalInfluenceToward}
\statementsource{Analysis-of-Boolean-Functions-by-Ryan-ODonnell}{
  pdf-ffbe9c7b3306-p016-b005,
  pdf-ffbe9c7b3306-p275-b001
}
\label{FABL.coalitionalInfluenceToward}
\leanok
\uses{FABL.BooleanFunction, FABL.Sign, FABL.coalitionSuccessProbability, FABL.uniformProbability}
Let $f:\{−1,1\}^n\to\{−1,1\}$, let $J$ be a set of coordinates, and let $b\in\{−1,1\}$.
The coalitional influence of $J$ toward $b$ is the probability that the coalition can
achieve output $b$, minus the probability that a uniformly random input already has
output $b$.
\end{definition}

\begin{definition}[Coalitional influence]
\lean{FABL.coalitionalInfluence}
\statementsource{Analysis-of-Boolean-Functions-by-Ryan-ODonnell}{
  pdf-ffbe9c7b3306-p016-b005,
  pdf-ffbe9c7b3306-p274-b007,
  pdf-ffbe9c7b3306-p275-b001
}
\label{FABL.coalitionalInfluence}
\leanok
\uses{FABL.BooleanFunction, FABL.CoalitionCanAttain, FABL.uniformProbability}
Let $f:\{−1,1\}^n\to\{−1,1\}$ and let $J$ be a set of coordinates. The coalitional
influence of $J$ is the probability, for a uniformly random input $x$, that the
coordinates in $J$ can be assigned values so that $f$ can attain either output $1$ or
output $-1$ while all coordinates outside $J$ remain as in $x$.
\end{definition}

\begin{theorem}[The function's value is attainable]
\lean{FABL.coalitionCanAttain_current}
\label{FABL.coalitionCanAttain_current}
\leanok
\uses{FABL.BooleanFunction, FABL.CoalitionCanAttain, FABL.SignCube}
\difficulty{1}
Let $n$ be a natural number, let $f\colon\{−1,1\}^n\to\{−1,1\}$, let
$J\subseteq\{1,\ldots,n\}$, and let $x\in\{−1,1\}^n$. Then $f(x)$ can be attained by an
input agreeing with $x$ at every coordinate outside $J$.
\end{theorem}

\begin{theorem}[Monotonicity of uniform probability]
\lean{FABL.uniformProbability_mono}
\label{FABL.uniformProbability_mono}
\leanok
\uses{FABL.uniformProbability}
\difficulty{4}
Let $\Omega$ be a finite nonempty set, and let $P$ and $Q$ be predicates on $\Omega$
such that every element satisfying $P$ also satisfies $Q$. Then the uniform probability
that $P$ holds is at most the uniform probability that $Q$ holds.
\end{theorem}

\begin{theorem}[Nonnegativity of Coalitional Influence]
\lean{FABL.coalitionalInfluenceToward_nonneg}
\label{FABL.coalitionalInfluenceToward_nonneg}
\leanok
\uses{FABL.BooleanFunction, FABL.CoalitionCanAttain, FABL.Sign, FABL.coalitionCanAttain_current, FABL.coalitionSuccessProbability, FABL.coalitionalInfluenceToward, FABL.uniformProbability_mono}
\difficulty{3}
Let $n$ be a nonnegative integer, let $f:\{-1,1\}^n\to\{-1,1\}$, let
$J\subseteq\{0,\ldots,n-1\}$, and let $b\in\{-1,1\}$. The coalitional influence
of $f$ toward $b$ is nonnegative.
\end{theorem}

\begin{theorem}[Coalitional influence lies between zero and one]
\lean{FABL.coalitionalInfluence_mem_Icc}
\label{FABL.coalitionalInfluence_mem_Icc}
\leanok
\uses{FABL.BooleanFunction, FABL.CoalitionCanAttain, FABL.SignCube, FABL.coalitionalInfluence, FABL.uniformProbability}
\difficulty{4}
Let $f:\{−1,1\}^n\to\{−1,1\}$ be a Boolean function and let $J$ be a set of coordinates.
The coalitional influence of $f$ on $J$ lies in $[0,1]$.
\end{theorem}

\begin{theorem}[Decomposition into directional influences]
\lean{FABL.coalitionalInfluence_eq_toward_add_toward}
\label{FABL.coalitionalInfluence_eq_toward_add_toward}
\leanok
\uses{FABL.BooleanFunction, FABL.CoalitionCanAttain, FABL.SignCube, FABL.coalitionCanAttain_current, FABL.coalitionSuccessProbability, FABL.coalitionalInfluence, FABL.coalitionalInfluenceToward, FABL.uniformProbability, FABL.uniformProbability_one_add_neg_one_eq_one}
\difficulty{5}
Let $f:\{−1,1\}^n\to\{−1,1\}$ be a Boolean function and let $J$ be a set of coordinates.
Its coalitional influence on $J$ equals the sum of its coalitional influences toward $1$
and toward $-1$.
\end{theorem}

\begin{theorem}[Attaining the Positive Sign by a Coalition]
\lean{FABL.coalitionCanAttain_one_iff_bribe_one}
\label{FABL.coalitionCanAttain_one_iff_bribe_one}
\leanok
\uses{FABL.BooleanFunction, FABL.CoalitionCanAttain, FABL.SignCube, FABL.bribeFunction, FABL.bribeInput, FABL.sign_le_one}
\difficulty{4}
Let $n$ be a nonnegative integer, let $f:\{-1,1\}^n\to\{-1,1\}$ be monotone, let
$J\subseteq\{1,\ldots,n\}$, and let $x\in\{-1,1\}^n$. There exists an input $y$ agreeing
with $x$ outside $J$ and satisfying $f(y)=1$ if and only if $f$ evaluated at the input
obtained from $x$ by setting every coordinate in $J$ to $1$ equals $1$.
\end{theorem}

\begin{theorem}[Attaining the Negative Sign by a Coalition]
\lean{FABL.coalitionCanAttain_neg_one_iff_bribe_neg_one}
\label{FABL.coalitionCanAttain_neg_one_iff_bribe_neg_one}
\leanok
\uses{FABL.BooleanFunction, FABL.CoalitionCanAttain, FABL.SignCube, FABL.bribeFunction, FABL.bribeInput, FABL.neg_one_le_sign}
\difficulty{4}
Let $n$ be a nonnegative integer, let $f:\{-1,1\}^n\to\{-1,1\}$ be monotone, let
$J\subseteq\{1,\ldots,n\}$, and let $x\in\{-1,1\}^n$. There exists an input $y$ agreeing
with $x$ outside $J$ and satisfying $f(y)=-1$ if and only if $f$ evaluated at the input
obtained from $x$ by setting every coordinate in $J$ to $-1$ equals $-1$.
\end{theorem}

\begin{theorem}[Positive Influence as a Mean Change]
\lean{FABL.coalitionalInfluenceToward_one_eq_mean_change}
\label{FABL.coalitionalInfluenceToward_one_eq_mean_change}
\leanok
\uses{FABL.BooleanFunction, FABL.BooleanFunction.toReal, FABL.CoalitionCanAttain, FABL.bribeFunction, FABL.coalitionCanAttain_one_iff_bribe_one, FABL.coalitionSuccessProbability, FABL.coalitionalInfluenceToward, FABL.mean, FABL.mean_eq_probability_one_sub_probability_neg_one, FABL.uniformProbability, FABL.uniformProbability_one_add_neg_one_eq_one}
\difficulty{4}
Let $n$ be a nonnegative integer, let $f:\{-1,1\}^n\to\{-1,1\}$ be monotone in the
coordinatewise order, and let $J\subseteq\{0,\ldots,n-1\}$. Write $\mathbb E f$ for
the uniform mean of the real-valued function obtained by viewing the signs output by $f$
as the numbers $\pm1$, and let $f_{J,+}$ be the function obtained by setting every
coordinate in $J$ to $+1$ before applying $f$. Then the coalitional influence of $f$
toward $+1$ equals $\bigl(\mathbb E f_{J,+}-\mathbb E f\bigr)/2$.
\end{theorem}

\begin{theorem}[Negative Influence as a Mean Change]
\lean{FABL.coalitionalInfluenceToward_neg_one_eq_mean_change}
\label{FABL.coalitionalInfluenceToward_neg_one_eq_mean_change}
\leanok
\uses{FABL.BooleanFunction, FABL.BooleanFunction.toReal, FABL.CoalitionCanAttain, FABL.bribeFunction, FABL.coalitionCanAttain_neg_one_iff_bribe_neg_one, FABL.coalitionSuccessProbability, FABL.coalitionalInfluenceToward, FABL.mean, FABL.mean_eq_probability_one_sub_probability_neg_one, FABL.uniformProbability, FABL.uniformProbability_one_add_neg_one_eq_one}
\difficulty{4}
Let $n$ be a nonnegative integer, let $f:\{-1,1\}^n\to\{-1,1\}$ be monotone in the
coordinatewise order, and let $J\subseteq\{0,\ldots,n-1\}$. Write $\mathbb E f$ for
the uniform mean of the real-valued function obtained by viewing the signs output by $f$
as the numbers $\pm1$. Let $f_{J,-}$ be the function obtained by setting every
coordinate in $J$ to $-1$ before applying $f$. Then the coalitional influence of $f$
toward $-1$ equals $\bigl(\mathbb E f-\mathbb E f_{J,-}\bigr)/2$.
\end{theorem}

\begin{theorem}[Attainment by a Singleton Coalition]
\lean{FABL.coalitionCanAttain_singleton_iff}
\label{FABL.coalitionCanAttain_singleton_iff}
\leanok
\uses{FABL.BooleanFunction, FABL.CoalitionCanAttain, FABL.Sign, FABL.SignCube, FABL.setCoordinate, FABL.setCoordinate_apply_of_ne}
\difficulty{4}
Let $n$ be a nonnegative integer, let $f:\{-1,1\}^n\to\{-1,1\}$, let $i$ be one of its
coordinates, let $b\in\{-1,1\}$, and let $x\in\{-1,1\}^n$. A coalition consisting only
of coordinate $i$ can make $f$ equal $b$ from $x$ if and only if at least one of the two
inputs obtained by setting coordinate $i$ to $1$ or to $-1$ has function value $b$.
\end{theorem}

\begin{theorem}[Singleton Coalition and Pivotality]
\lean{FABL.coalitionCanAttain_singleton_both_iff_isPivotal}
\label{FABL.coalitionCanAttain_singleton_both_iff_isPivotal}
\leanok
\uses{FABL.BooleanFunction, FABL.CoalitionCanAttain, FABL.IsPivotal, FABL.SignCube, FABL.coalitionCanAttain_singleton_iff, FABL.isPivotal_iff_setCoordinate_ne, FABL.setCoordinate}
\difficulty{4}
Let $n$ be a nonnegative integer, let $f:\{-1,1\}^n\to\{-1,1\}$, let $i$ be one of its
coordinates, and let $x\in\{-1,1\}^n$. A coalition consisting only of coordinate $i$ can
attain both output signs from $x$ if and only if flipping coordinate $i$ changes the
value of $f$ at $x$.
\end{theorem}

\begin{theorem}[Singleton coalitional influence]
\lean{FABL.coalitionalInfluence_singleton}
\label{FABL.coalitionalInfluence_singleton}
\leanok
\uses{FABL.BooleanFunction, FABL.CoalitionCanAttain, FABL.booleanInfluence, FABL.coalitionCanAttain_singleton_both_iff_isPivotal, FABL.coalitionalInfluence, FABL.uniformProbability}
\difficulty{3}
Let $n$ be a nonnegative integer, let $f:\{-1,1\}^n\to\{-1,1\}$, and let $i$ be one of
the $n$ coordinates. The coalitional influence of $f$ with respect to the singleton set
$\{i\}$ equals its Boolean influence at coordinate $i$, namely the probability that
flipping coordinate $i$ changes the value of $f$.
\end{theorem}

\begin{theorem}[Singleton coalitional influence]
\lean{FABL.coalitionalInfluenceToward_singleton}
\label{FABL.coalitionalInfluenceToward_singleton}
\leanok
\uses{FABL.BooleanFunction, FABL.CoalitionCanAttain, FABL.IsPivotal, FABL.Sign, FABL.SignCube, FABL.booleanInfluence, FABL.coalitionCanAttain_singleton_iff, FABL.coalitionSuccessProbability, FABL.coalitionalInfluenceToward, FABL.coordinateExpectation_apply, FABL.isPivotal_iff_setCoordinate_ne, FABL.mean, FABL.mean_coordinateExpectation, FABL.setCoordinate, FABL.uniformProbability}
\difficulty{6}
Let $f:\{−1,1\}^n\to\{−1,1\}$ be a Boolean function, let $i$ be a coordinate, and let
$b\in\{−1,1\}$. The influence of the singleton coalition $\{i\}$ toward $b$ equals half
the probability that flipping coordinate $i$ changes the value of $f$.
\end{theorem}

\begin{theorem}[Singleton influence identities]
\lean{FABL.exercise9_25_singleton_equalities}
\label{FABL.exercise9_25_singleton_equalities}
\leanok
\uses{FABL.BooleanFunction, FABL.booleanInfluence, FABL.coalitionalInfluence, FABL.coalitionalInfluenceToward, FABL.coalitionalInfluenceToward_singleton, FABL.coalitionalInfluence_singleton}
\difficulty{3}
Let $f:\{-1,1\}^n\to\{-1,1\}$ and let $i$ be a coordinate. The uniform probability that
flipping coordinate $i$ changes $f$ equals the coalitional influence of the singleton
$\{i\}$, and this influence is twice the coalitional influence toward either output:
$\operatorname{Inf}_i(f)=\operatorname{Inf}_{\{i\}}(f)=2\operatorname{Inf}_{\{i\},+1}(f)=2\operatorname{Inf}_{\{i\},-1}(f)$.
\end{theorem}

\begin{theorem}[Parity has maximal coalitional influence]
\lean{FABL.coalitionalInfluence_parity_univ}
\label{FABL.coalitionalInfluence_parity_univ}
\leanok
\uses{FABL.CoalitionCanAttain, FABL.coalitionalInfluence, FABL.flipCoordinate, FABL.parityFunction, FABL.parityFunction_univ_flipCoordinate, FABL.setCoordinate_apply_of_ne, FABL.uniformProbability}
\difficulty{5}
Let $n$ be a nonnegative integer and let $J$ be a nonempty set of coordinates among $n$
coordinates. The coalitional influence of the parity function on all $n$ coordinates,
with respect to $J$, equals $1$.
\end{theorem}

\begin{theorem}[Monotonicity of the mean]
\lean{FABL.mean_toReal_mono}
\label{FABL.mean_toReal_mono}
\leanok
\uses{FABL.BooleanFunction, FABL.BooleanFunction.toReal, FABL.mean, FABL.monotone_signValue}
\difficulty{3}
Let $n$ be a nonnegative integer, and let $f$ and $g$ map the sign cube to the two
signs. If $f(x)\le g(x)$ for every input $x$, then the mean of their real-valued sign
representations satisfies $\mathbb E[f]\le\mathbb E[g]$.
\end{theorem}

\begin{theorem}[Monotonicity under positive bribery]
\lean{FABL.mean_bribeFunction_mono_set_one}
\label{FABL.mean_bribeFunction_mono_set_one}
\leanok
\uses{FABL.BooleanFunction, FABL.BooleanFunction.toReal, FABL.bribeFunction, FABL.bribeInput, FABL.mean, FABL.mean_toReal_mono, FABL.sign_le_one}
\difficulty{4}
Let $f$ be a monotone Boolean function on the $n$-dimensional sign cube, and let
$J\subseteq K$ be sets of coordinates. The mean after replacing the coordinates in $J$
by $1$ is at most the mean after replacing the coordinates in $K$ by $1$.
\end{theorem}

\begin{theorem}[Monotonicity under negative bribery]
\lean{FABL.mean_bribeFunction_mono_set_neg_one}
\label{FABL.mean_bribeFunction_mono_set_neg_one}
\leanok
\uses{FABL.BooleanFunction, FABL.BooleanFunction.toReal, FABL.bribeFunction, FABL.bribeInput, FABL.mean, FABL.mean_toReal_mono, FABL.neg_one_le_sign}
\difficulty{4}
Let $f$ be a monotone Boolean function on the $n$-dimensional sign cube, and let
$J\subseteq K$ be sets of coordinates. The mean after replacing the coordinates in $K$
by $-1$ is at most the mean after replacing the coordinates in $J$ by $-1$.
\end{theorem}

\begin{theorem}[A fixed coalition with high influence]
\lean{FABL.proposition9_27_fixed_coalition}
\proofsource{Analysis-of-Boolean-Functions-by-Ryan-ODonnell}{
  pdf-ffbe9c7b3306-p265-b006,
  pdf-ffbe9c7b3306-p266-b001
}
\label{FABL.proposition9_27_fixed_coalition}
\leanok
\uses{FABL.BooleanFunction, FABL.BooleanFunction.toReal, FABL.bribeFunction, FABL.coalitionalInfluence, FABL.coalitionalInfluenceToward_neg_one_eq_mean_change, FABL.coalitionalInfluenceToward_one_eq_mean_change, FABL.coalitionalInfluence_eq_toward_add_toward, FABL.mean, FABL.mean_bribeFunction_mono_set_neg_one, FABL.mean_bribeFunction_mono_set_one, FABL.proposition9_27_negative, FABL.proposition9_27_positive}
\difficulty{5}
Let $n\ge 2$, and let $f:\{-1,1\}^n\to\{-1,1\}$ be monotone, with uniform mean between
$-99/100$ and $99/100$, inclusive. There is a set $J$ of coordinates, of size at most
$20000n/\log n$, such that for at least $99/100$ of uniformly chosen inputs, the
coalition controlling $J$ can make the output either sign by changing coordinates in $J$
while leaving all other coordinates fixed.
\end{theorem}

\begin{definition}[Positive bias]
\lean{FABL.positiveBias}
\label{FABL.positiveBias}
\leanok
\uses{FABL.BooleanFunction, FABL.uniformProbability}
For a Boolean function $f$ on the $n$-dimensional sign cube, its positive bias is the
uniform probability that $f(x)=1$.
\end{definition}

\begin{theorem}[Positive bias from the mean]
\lean{FABL.positiveBias_eq_mean}
\label{FABL.positiveBias_eq_mean}
\leanok
\uses{FABL.BooleanFunction, FABL.BooleanFunction.toReal, FABL.mean, FABL.mean_eq_probability_one_sub_probability_neg_one, FABL.positiveBias, FABL.uniformProbability_one_add_neg_one_eq_one}
\difficulty{3}
For every Boolean function $f$ on a finite sign cube, its positive bias equals
$(1+\mathbb{E}[f])/2$, where the mean is taken after interpreting the output sign as the
real number $+1$ or $-1$.
\end{theorem}

\begin{theorem}[Bounds on positive bias]
\lean{FABL.positiveBias_mem_Icc}
\label{FABL.positiveBias_mem_Icc}
\leanok
\uses{FABL.BooleanFunction, FABL.mean_toReal_mem_Icc, FABL.positiveBias, FABL.positiveBias_eq_mean}
\difficulty{3}
For every Boolean function on a finite sign cube, its positive bias lies in the interval
$[0,1]$.
\end{theorem}

\begin{theorem}[Variance from positive bias]
\lean{FABL.variance_eq_four_mul_positiveBias}
\label{FABL.variance_eq_four_mul_positiveBias}
\leanok
\uses{FABL.BooleanFunction, FABL.BooleanFunction.toReal, FABL.positiveBias, FABL.uniformProbability_one_add_neg_one_eq_one, FABL.variance, FABL.variance_eq_four_mul_probabilities}
\difficulty{3}
Let $f$ be a Boolean function on the $n$-dimensional sign cube, and let $p$ be the
uniform probability that $f$ takes the value $+1$. Then $\operatorname{Var}(f)=4p(1-p)$.
\end{theorem}

\begin{definition}[Bribery rate]
\lean{FABL.briberyRate}
\statementsource{Analysis-of-Boolean-Functions-by-Ryan-ODonnell}{pdf-ffbe9c7b3306-p266-b001}
\label{FABL.briberyRate}
\leanok
For a natural number $n$, define the bribery rate to be $\log(n)/(100n)$, where $\log$
is the natural logarithm.
\end{definition}

\begin{theorem}[Positivity of the bribery rate]
\lean{FABL.briberyRate_pos}
\label{FABL.briberyRate_pos}
\leanok
\uses{FABL.briberyRate}
\difficulty{3}
For every integer $n\ge 2$, the quantity $\frac{\log n}{100n}$ is positive, with the
logarithm and arithmetic interpreted over the real numbers.
\end{theorem}

\begin{theorem}[Upper bound for the bribery rate]
\lean{FABL.briberyRate_le_one_half}
\label{FABL.briberyRate_le_one_half}
\leanok
\uses{FABL.briberyRate}
\difficulty{3}
For every integer $n\ge 2$, the quantity $\frac{\log n}{100n}$ is at most $\frac12$,
with the logarithm and arithmetic interpreted over the real numbers.
\end{theorem}

\begin{theorem}[Maximum influence coordinate positive bias step]
\lean{FABL.maximumInfluenceCoordinate_positiveBias_step}
\statementsource{Analysis-of-Boolean-Functions-by-Ryan-ODonnell}{pdf-ffbe9c7b3306-p266-b001}
\label{FABL.maximumInfluenceCoordinate_positiveBias_step}
\leanok
\uses{FABL.BooleanFunction, FABL.BooleanFunction.toReal, FABL.booleanInfluence, FABL.booleanInfluence_eq_influence_toReal, FABL.booleanInfluence_maximumInfluenceCoordinate, FABL.briberyRate, FABL.influence_eq_fourierCoeff_singleton_of_monotone, FABL.kkl, FABL.maximumInfluence, FABL.maximumInfluenceCoordinate, FABL.mean, FABL.mean_fixCoordinate_one, FABL.positiveBias, FABL.positiveBias_eq_mean, FABL.setCoordinate, FABL.variance, FABL.variance_eq_four_mul_positiveBias}
\difficulty{5}
% dataset-pipeline: stub -- prose pending build_informal_statements.py
Statement pending regeneration.
\end{theorem}

\begin{definition}[Bribery sets up to a target]
\lean{FABL.briberySetsTo}
\statementsource{Analysis-of-Boolean-Functions-by-Ryan-ODonnell}{pdf-ffbe9c7b3306-p265-b007}
\label{FABL.briberySetsTo}
\leanok
\uses{FABL.BooleanFunction, FABL.BooleanFunction.toReal, FABL.bribeFunction, FABL.maximumInfluenceCoordinate, FABL.mean}
Let $n>0$, let $f:\{−1,1\}^n\to\{−1,1\}$, and let $r\in\mathbb R$. Define sets of
coordinates recursively by $J_0=\varnothing$. Given $J_t$, form the function obtained by
setting every coordinate in $J_t$ to $1$ before evaluating $f$. If its uniform mean is
less than $r$, add a coordinate of maximum influence for that modified function to
obtain $J_{t+1}$; otherwise keep $J_{t+1}=J_t$.
\end{definition}

\begin{theorem}[One step of the bribery procedure]
\lean{FABL.briberySetsTo_succ}
\label{FABL.briberySetsTo_succ}
\leanok
\uses{FABL.BooleanFunction, FABL.BooleanFunction.toReal, FABL.bribeFunction, FABL.briberySetsTo, FABL.maximumInfluenceCoordinate, FABL.mean}
\difficulty{1}
Let $n$ be a positive integer, let $f$ be a Boolean function on sign vectors of length
$n$, let $q$ be a real target, and let $t$ be a nonnegative integer. The set at step
$t+1$ is obtained by adding to the set at step $t$ a coordinate of maximum influence of
the function bribed on that set if its uniform mean is below $q$; if its mean is at
least $q$, the set is unchanged.
\end{theorem}

\begin{theorem}[Size of the bribery set after finitely many steps]
\lean{FABL.card_briberySetsTo_le}
\label{FABL.card_briberySetsTo_le}
\leanok
\uses{FABL.BooleanFunction, FABL.briberySetsTo, FABL.briberySetsTo_succ}
\difficulty{4}
Let $n>0$, let $f:\{-1,1\}^n\to\{-1,1\}$, let $a\in\mathbb R$, and let
$t\in\mathbb N$. Starting from the empty set of coordinates, perform $t$ steps: at
each step, set every coordinate already chosen to $1$, and if the resulting function has
mean less than $a$, add a coordinate of maximum influence; otherwise leave the set
unchanged. The set of chosen coordinates after $t$ steps has cardinality at most $t$.
\end{theorem}

\begin{theorem}[Stability of the bribery set]
\lean{FABL.briberySetsTo_stable}
\proofsource{Analysis-of-Boolean-Functions-by-Ryan-ODonnell}{pdf-ffbe9c7b3306-p266-b001}
\label{FABL.briberySetsTo_stable}
\leanok
\uses{FABL.BooleanFunction, FABL.BooleanFunction.toReal, FABL.bribeFunction, FABL.briberySetsTo, FABL.briberySetsTo_succ, FABL.mean}
\difficulty{4}
Let $n$ be a positive integer, let $f$ be a Boolean function on sign vectors of length
$n$, and let $q$ be a real target. For a nonnegative integer $t$, let $J_t$ be the set
obtained by starting with the empty set and, at each step, adding a coordinate of
maximum influence if the uniform mean of $f$ after replacing the coordinates in the
current set by $1$ is below $q$, and otherwise keeping the set unchanged. If that mean
at step $t$ is at least $q$, then $J_{t+s}=J_t$ for every nonnegative integer $s$.
\end{theorem}

\begin{theorem}[Fresh coordinate of maximum influence]
\lean{FABL.maximumInfluenceCoordinate_fresh_of_bias}
\proofsource{Analysis-of-Boolean-Functions-by-Ryan-ODonnell}{pdf-ffbe9c7b3306-p266-b001}
\label{FABL.maximumInfluenceCoordinate_fresh_of_bias}
\leanok
\uses{FABL.BooleanFunction, FABL.BooleanFunction.toReal, FABL.booleanInfluence, FABL.booleanInfluence_bribeFunction_eq_zero_of_mem, FABL.booleanInfluence_maximumInfluenceCoordinate, FABL.bribeFunction, FABL.briberyRate, FABL.briberyRate_pos, FABL.kkl, FABL.maximumInfluence, FABL.maximumInfluenceCoordinate, FABL.positiveBias, FABL.variance, FABL.variance_eq_four_mul_positiveBias}
\difficulty{4}
Let $n\ge 2$, let $f:\{-1,1\}^n\to\{-1,1\}$ be a Boolean function, and let
$J\subseteq\{1,ldots,n\}$. Form a new function by setting every coordinate in $J$ to $1$
before evaluating $f$. If the fraction of inputs on which this new function equals $1$
is strictly between $0$ and $1$, then a coordinate of maximum influence of the new
function lies outside $J$.
\end{theorem}

\begin{theorem}[Size of the Recursive Bribery Set]
\lean{FABL.card_briberySetsTo_eq_of_mean_lt}
\proofsource{Analysis-of-Boolean-Functions-by-Ryan-ODonnell}{pdf-ffbe9c7b3306-p266-b001}
\label{FABL.card_briberySetsTo_eq_of_mean_lt}
\leanok
\uses{FABL.BooleanFunction, FABL.BooleanFunction.toReal, FABL.bribeFunction, FABL.bribeFunction_empty, FABL.briberySetsTo, FABL.briberySetsTo_succ, FABL.maximumInfluenceCoordinate, FABL.maximumInfluenceCoordinate_fresh_of_bias, FABL.mean, FABL.mean_bribeFunction_mono_set_one, FABL.positiveBias, FABL.positiveBias_eq_mean}
\difficulty{5}
Let $n\ge 2$, and let $f$ be a monotone Boolean function on the sign cube of dimension
$n$ with positive bias. Let $a<1$ be a real target and let $t\le n$. Starting from the
empty set, form the recursively selected bribery set by adjoining a maximum-influence
coordinate whenever setting the coordinates already selected to $+1$ leaves the mean of
the resulting function below $a$, and otherwise leaving the set unchanged. If, after $t$
steps, that mean is still below $a$, then the selected set has cardinality exactly $t$.
\end{theorem}

\begin{theorem}[Odds step for positive bias]
\lean{FABL.positiveBias_odds_step}
\label{FABL.positiveBias_odds_step}
\leanok
\uses{FABL.BooleanFunction, FABL.briberyRate, FABL.briberyRate_pos, FABL.maximumInfluenceCoordinate, FABL.maximumInfluenceCoordinate_positiveBias_step, FABL.positiveBias, FABL.positiveBias_mem_Icc, FABL.setCoordinate}
\difficulty{4}
Let $n\ge 2$, and let $f$ be a monotone Boolean function on the $n$-dimensional sign
cube whose positive bias $p$ satisfies $0<p<1$. Choose a coordinate $i$ of maximum
influence, where influence is the fraction of inputs on which flipping that coordinate
changes the function, and let $p'$ be the fraction of inputs on which $f$ is positive
after coordinate $i$ is set to $+1$. Then, with $r=\log(n)/(100n)$, $(1+2r)p(1-p')\le
p'(1-p)$.
\end{theorem}

\begin{theorem}[Bribery-set odds invariant]
\lean{FABL.briberySetsTo_odds_invariant}
\label{FABL.briberySetsTo_odds_invariant}
\leanok
\uses{FABL.BooleanFunction, FABL.bribeFunction, FABL.bribeFunction_empty, FABL.bribeFunction_insert, FABL.briberyRate, FABL.briberyRate_pos, FABL.briberySetsTo, FABL.briberySetsTo_succ, FABL.maximumInfluenceCoordinate, FABL.maximumInfluenceCoordinate_fresh_of_bias, FABL.mean_bribeFunction_mono_set_one, FABL.monotone_bribeFunction, FABL.positiveBias, FABL.positiveBias_eq_mean, FABL.positiveBias_mem_Icc, FABL.positiveBias_odds_step, FABL.setCoordinate}
\difficulty{6}
Let $n\ge 2$, and let $f:\{-1,1\}^n\to\{-1,1\}$ be monotone with positive bias
$p=\Pr[f(x)=1]>0$ under the uniform distribution. Fix a real target $q<1$. Starting from
$J_0=\varnothing$, define $J_{t+1}$ by leaving $J_t$ unchanged if the uniform mean of
$x\mapsto f(x^{J_t})$ is at least $q$, and otherwise adjoining to $J_t$ a coordinate of
maximum influence for that function; here $x^{J_t}$ is obtained by setting every
coordinate in $J_t$ to $1$. For every integer $t\ge0$, if $p_t=\Pr[f(x^{J_t})=1]$ and
$r=\frac{\log n}{100n}$, then $(1+2r)^{|J_t|}p(1-p_t)\le p_t(1-p)$.
\end{theorem}

\begin{theorem}[Mean bribe function univ one of positive bias]
\lean{FABL.mean_bribeFunction_univ_one_of_positiveBias}
\statementsource{Analysis-of-Boolean-Functions-by-Ryan-ODonnell}{pdf-ffbe9c7b3306-p266-b001}
\label{FABL.mean_bribeFunction_univ_one_of_positiveBias}
\leanok
\uses{FABL.BooleanFunction, FABL.BooleanFunction.toReal, FABL.Sign, FABL.bribeFunction, FABL.bribeInput, FABL.mean, FABL.neg_one_le_sign, FABL.positiveBias, FABL.sign_le_one, FABL.uniformProbability}
\difficulty{4}
% dataset-pipeline: stub -- prose pending build_informal_statements.py
Statement pending regeneration.
\end{theorem}

\begin{theorem}[Mean bribery sets to final ge target]
\lean{FABL.mean_briberySetsTo_final_ge_target}
\statementsource{Analysis-of-Boolean-Functions-by-Ryan-ODonnell}{pdf-ffbe9c7b3306-p266-b001}
\label{FABL.mean_briberySetsTo_final_ge_target}
\leanok
\uses{FABL.BooleanFunction, FABL.BooleanFunction.toReal, FABL.bribeFunction, FABL.briberySetsTo, FABL.card_briberySetsTo_eq_of_mean_lt, FABL.mean, FABL.mean_bribeFunction_univ_one_of_positiveBias, FABL.positiveBias}
\difficulty{4}
% dataset-pipeline: stub -- prose pending build_informal_statements.py
Statement pending regeneration.
\end{theorem}

\begin{theorem}[Mean below target during bribery]
\lean{FABL.mean_lt_target_of_final_card_gt}
\label{FABL.mean_lt_target_of_final_card_gt}
\leanok
\uses{FABL.BooleanFunction, FABL.BooleanFunction.toReal, FABL.bribeFunction, FABL.briberySetsTo, FABL.briberySetsTo_stable, FABL.card_briberySetsTo_le, FABL.mean}
\difficulty{4}
Let $n>0$, let $f:\{−1,1\}^n\to\{−1,1\}$, and let $\tau\in\mathbb R$. Starting with
$J_0=\varnothing$, construct $J_{k+1}$ from $J_k$ by adding a coordinate of maximum
influence for the function obtained from $f$ by setting all coordinates in $J_k$ to $1$
if that function's mean is below $\tau$; otherwise leave the set unchanged. If $t\le n$
and the set $J_n$ has cardinality greater than $t$, then the mean of the function
obtained by setting the coordinates in $J_t$ to $1$ is strictly less than $\tau$.
\end{theorem}

\begin{theorem}[A logarithmic bound for the bribery rate]
\lean{FABL.log_one_add_two_mul_briberyRate_ge}
\label{FABL.log_one_add_two_mul_briberyRate_ge}
\leanok
\uses{FABL.briberyRate, FABL.briberyRate_le_one_half, FABL.briberyRate_pos}
\difficulty{4}
For every integer $n \ge 2$, let $r=\log(n)/(100n)$. Then $r \le \log(1+2r)$.
\end{theorem}

\begin{theorem}[Positive bribery bound]
\lean{FABL.exercise9_27_positive}
\proofsource{Analysis-of-Boolean-Functions-by-Ryan-ODonnell}{pdf-ffbe9c7b3306-p275-b005}
\label{FABL.exercise9_27_positive}
\leanok
\uses{FABL.BooleanFunction, FABL.BooleanFunction.toReal, FABL.bribeFunction, FABL.briberyRate, FABL.briberyRate_pos, FABL.briberySetsTo, FABL.briberySetsTo_odds_invariant, FABL.card_briberySetsTo_eq_of_mean_lt, FABL.log_one_add_two_mul_briberyRate_ge, FABL.mean, FABL.mean_briberySetsTo_final_ge_target, FABL.mean_lt_target_of_final_card_gt, FABL.positiveBias, FABL.positiveBias_eq_mean, FABL.positiveBias_mem_Icc, FABL.variance, FABL.variance_eq_four_mul_positiveBias, FABL.variance_eq_four_mul_probabilities}
\difficulty{6}
Let $f:\{-1,1\}^n\to\{-1,1\}$ be monotone, with $2\le n$. Suppose $\delta>0$ and
$\delta$ is at most the variance of the real-valued sign representation of $f$. Let
$0<\varepsilon<1/2$. Then there is a set $J$ of coordinates such that setting every
coordinate in $J$ to $+1$ gives a function with mean at least $1-\varepsilon$, and
$|J|\le 1+\frac{100n}{\log n}\log\!\left(\frac{8}{\varepsilon\delta}\right)$.
\end{theorem}

\begin{theorem}[Variance under Boolean duality]
\lean{FABL.variance_dualBooleanFunction}
\label{FABL.variance_dualBooleanFunction}
\leanok
\uses{FABL.BooleanFunction, FABL.BooleanFunction.toReal, FABL.dualBooleanFunction, FABL.mean_dualBooleanFunction, FABL.variance, FABL.variance_eq_four_mul_probabilities}
\difficulty{3}
Let $f$ be a Boolean function on the $n$-dimensional sign cube, and define its dual by
$f^*(x)=-f(-x)$. Then $\operatorname{Var}(f^*)=\operatorname{Var}(f)$.
\end{theorem}

\begin{theorem}[Exercise9 27 coalition]
\lean{FABL.exercise9_27_coalition}
\statementsource{Analysis-of-Boolean-Functions-by-Ryan-ODonnell}{pdf-ffbe9c7b3306-p275-b005}
\label{FABL.exercise9_27_coalition}
\leanok
\uses{FABL.BooleanFunction, FABL.BooleanFunction.toReal, FABL.bribeFunction, FABL.coalitionalInfluence, FABL.coalitionalInfluenceToward_neg_one_eq_mean_change, FABL.coalitionalInfluenceToward_one_eq_mean_change, FABL.coalitionalInfluence_eq_toward_add_toward, FABL.dualBooleanFunction, FABL.exercise9_27_positive, FABL.mean, FABL.mean_bribeFunction_dual_one, FABL.mean_bribeFunction_mono_set_neg_one, FABL.mean_bribeFunction_mono_set_one, FABL.monotone_dualBooleanFunction, FABL.variance, FABL.variance_dualBooleanFunction}
\difficulty{5}
% dataset-pipeline: stub -- prose pending build_informal_statements.py
Statement pending regeneration.
\end{theorem}

\begin{definition}[Coordinate polarization]
\lean{FABL.coordinatePolarization}
\statementsource{Analysis-of-Boolean-Functions-by-Ryan-ODonnell}{pdf-ffbe9c7b3306-p065-b005}
\label{FABL.coordinatePolarization}
\leanok
\uses{FABL.BooleanFunction, FABL.setCoordinate}
Let $f:\{−1,1\}^n\to\{−1,1\}$ and let $i$ be a coordinate. Define a new Boolean function
by assigning, at each input $x$, the larger of $f(x^{i,+})$ and $f(x^{i,-})$ when
$x_i=1$, and the smaller when $x_i=-1$, where $x^{i,+}$ and $x^{i,-}$ are obtained from
$x$ by setting coordinate $i$ to $1$ and $-1$, respectively.
\end{definition}

\begin{theorem}[Polarization at the positive coordinate value]
\lean{FABL.coordinatePolarization_setCoordinate_one}
\label{FABL.coordinatePolarization_setCoordinate_one}
\leanok
\uses{FABL.BooleanFunction, FABL.SignCube, FABL.coordinatePolarization, FABL.setCoordinate}
\difficulty{1}
Let $n$ be a nonnegative integer, let $f$ be a function from the $n$-dimensional sign
cube $\{-1,1\}^n$ to $\{-1,1\}$, and let $i$ be a coordinate and $x$ a point of the
sign cube. The coordinate polarization of $f$, evaluated at the point obtained from $x$
by setting coordinate $i$ to $1$, equals the maximum of the values of $f$ at the two
points obtained by setting coordinate $i$ to $1$ and to $-1$.
\end{theorem}

\begin{theorem}[Polarization at the negative coordinate]
\lean{FABL.coordinatePolarization_setCoordinate_neg_one}
\label{FABL.coordinatePolarization_setCoordinate_neg_one}
\leanok
\uses{FABL.BooleanFunction, FABL.SignCube, FABL.coordinatePolarization, FABL.setCoordinate}
\difficulty{1}
Let $f:\{−1,1\}^n\to\{−1,1\}$, let $i\in\{1,\ldots,n\}$, and let $x\in\{−1,1\}^n$. The
coordinate polarization of $f$ in coordinate $i$, evaluated at the vector obtained from
$x$ by setting coordinate $i$ to $-1$, equals the minimum of the values of $f$ on the
two vectors obtained from $x$ by setting coordinate $i$ to $1$ and to $-1$,
respectively.
\end{theorem}

\begin{theorem}[Coordinate expectation is preserved by polarization]
\lean{FABL.coordinateExpectation_coordinatePolarization_toReal}
\label{FABL.coordinateExpectation_coordinatePolarization_toReal}
\leanok
\uses{FABL.BooleanFunction, FABL.BooleanFunction.toReal, FABL.coordinateExpectation, FABL.coordinateExpectation_apply, FABL.coordinatePolarization, FABL.coordinatePolarization_setCoordinate_neg_one, FABL.coordinatePolarization_setCoordinate_one, FABL.setCoordinate}
\difficulty{4}
Let $n\in\mathbb N$, let $f:\{−1,1\}^n\to\{−1,1\}$, and choose a coordinate
$i\in\mathrm{Fin}(n)$. The coordinate expectation of the real-valued polarization of $f$
in coordinate $i$ equals the coordinate expectation of $f$ after interpreting its signs
as real numbers.
\end{theorem}

\begin{theorem}[Mean under coordinate polarization]
\lean{FABL.mean_coordinatePolarization}
\label{FABL.mean_coordinatePolarization}
\leanok
\uses{FABL.BooleanFunction, FABL.BooleanFunction.toReal, FABL.coordinateExpectation, FABL.coordinateExpectation_coordinatePolarization_toReal, FABL.coordinatePolarization, FABL.mean, FABL.mean_coordinateExpectation}
\difficulty{3}
Let $f$ be a Boolean function on the sign cube of dimension $n$, and choose any
coordinate $i$. Form a new function that, at inputs with coordinate $i$ equal to $+1$,
takes the maximum of the values of $f$ with coordinate $i$ set to either sign, and at
inputs with coordinate $i$ equal to $-1$, takes the minimum of those two values. The
resulting function has the same mean as $f$.
\end{theorem}

\begin{theorem}[Polarization preserves output probabilities]
\lean{FABL.uniformProbability_coordinatePolarization_eq}
\label{FABL.uniformProbability_coordinatePolarization_eq}
\leanok
\uses{FABL.BooleanFunction, FABL.Sign, FABL.coordinatePolarization, FABL.mean_coordinatePolarization, FABL.mean_eq_probability_one_sub_probability_neg_one, FABL.uniformProbability, FABL.uniformProbability_one_add_neg_one_eq_one}
\difficulty{4}
Let $n$ be a nonnegative integer, let $f$ be a sign-valued function on $\{-1,1\}^n$,
let $i$ be a coordinate, and let $b\in\{-1,1\}$. The uniform probability that the
coordinate polarization of $f$ at $i$ equals $b$ is the same as the uniform probability
that $f$ equals $b$.
\end{theorem}

\begin{theorem}[Commuting distinct coordinate updates]
\lean{FABL.setCoordinate_comm_of_ne}
\label{FABL.setCoordinate_comm_of_ne}
\leanok
\uses{FABL.Sign, FABL.SignCube, FABL.setCoordinate}
\difficulty{3}
Let $n$ be a natural number, let $x$ be a sign vector on $n$ coordinates, and let $i$
and $j$ be distinct coordinates. Replacing coordinate $i$ by $a$ and coordinate $j$ by
$b$ gives the same vector as performing these two replacements in the opposite order,
for any signs $a$ and $b$.
\end{theorem}

\begin{theorem}[Coordinate monotonicity after polarization]
\lean{FABL.isMonotoneInCoordinate_coordinatePolarization_self}
\proofsource{Analysis-of-Boolean-Functions-by-Ryan-ODonnell}{pdf-ffbe9c7b3306-p065-b005}
\label{FABL.isMonotoneInCoordinate_coordinatePolarization_self}
\leanok
\uses{FABL.BooleanFunction, FABL.IsMonotoneInCoordinate, FABL.coordinatePolarization, FABL.coordinatePolarization_setCoordinate_neg_one, FABL.coordinatePolarization_setCoordinate_one}
\difficulty{1}
Let $n$ be a nonnegative integer, let $f$ be a function from the $n$-dimensional sign
cube to $\{-1,1\}$, and let $i$ be a coordinate. The coordinate polarization of $f$ in
coordinate $i$ is monotone in coordinate $i$: for every input $x$, its value at the
point obtained from $x$ by setting coordinate $i$ to $-1$ is at most its value at the
point obtained by setting coordinate $i$ to $1$.
\end{theorem}

\begin{theorem}[Polarization preserves coordinate monotonicity]
\lean{FABL.isMonotoneInCoordinate_coordinatePolarization}
\proofsource{Analysis-of-Boolean-Functions-by-Ryan-ODonnell}{pdf-ffbe9c7b3306-p065-b005}
\label{FABL.isMonotoneInCoordinate_coordinatePolarization}
\leanok
\uses{FABL.BooleanFunction, FABL.IsMonotoneInCoordinate, FABL.coordinatePolarization, FABL.isMonotoneInCoordinate_coordinatePolarization_self, FABL.setCoordinate, FABL.setCoordinate_apply_of_ne, FABL.setCoordinate_comm_of_ne}
\difficulty{5}
Let $n$ be a nonnegative integer, let $f$ be a sign-valued function on the
$n$-dimensional Boolean cube, and let $i,j$ be coordinates. If $f$ is monotone in
coordinate $j$, then its polarization in coordinate $i$ is also monotone in coordinate
$j$.
\end{theorem}

\begin{theorem}[Attainability under coordinate polarization]
\lean{FABL.coalitionCanAttain_coordinatePolarization_imp_of_mem}
\label{FABL.coalitionCanAttain_coordinatePolarization_imp_of_mem}
\leanok
\uses{FABL.BooleanFunction, FABL.CoalitionCanAttain, FABL.Sign, FABL.SignCube, FABL.coordinatePolarization, FABL.setCoordinate, FABL.setCoordinate_apply_of_ne}
\difficulty{5}
Let $n$ be a natural number, let $f\colon\{−1,1\}^n\to\{−1,1\}$, let
$J\subseteq\{1,\ldots,n\}$, let $b\in\{−1,1\}$, and choose $i\in J$ and
$x\in\{−1,1\}^n$. Define the coordinate-polarized function $f_i^{\mathrm{pol}}$ by
taking, at an input $z$, the maximum of $f$ at the two inputs obtained by setting
coordinate $i$ to $1$ and $-1$ if $z_i=1$, and their minimum if $z_i=-1$. If one can
obtain output $b$ from $f_i^{\mathrm{pol}}$ at some input agreeing with $x$ outside $J$,
then one can obtain output $b$ from $f$ at some input agreeing with $x$ outside $J$.
\end{theorem}

\begin{theorem}[Attainability at both coordinate settings]
\lean{FABL.coalitionCanAttain_coordinatePolarization_one_set_one_imp}
\label{FABL.coalitionCanAttain_coordinatePolarization_one_set_one_imp}
\leanok
\uses{FABL.BooleanFunction, FABL.CoalitionCanAttain, FABL.Sign, FABL.SignCube, FABL.coordinatePolarization, FABL.setCoordinate, FABL.setCoordinate_apply_of_ne}
\difficulty{4}
Let $n$ be a natural number, let $f\colon\{−1,1\}^n\to\{−1,1\}$, let
$J\subseteq\{1,\ldots,n\}$, and choose a coordinate $i\notin J$ and an input
$x\in\{−1,1\}^n$. Define the coordinate-polarized function $f_i^{\mathrm{pol}}$ by
taking, at an input $z$, the maximum of $f$ at the two inputs obtained by setting
coordinate $i$ to $1$ and $-1$ if $z_i=1$, and their minimum if $z_i=-1$. If
$f_i^{\mathrm{pol}}$ can attain $1$ at an input agreeing with $x$ with coordinate $i$
set to $1$, and agreeing with that input outside $J$, then $f$ can attain $1$ from an
input with coordinate $i$ set to $1$ or from one with coordinate $i$ set to $-1$,
allowing changes only in $J$ from the respective input.
\end{theorem}

\begin{theorem}[Attainability of plus one after polarization]
\lean{FABL.coalitionCanAttain_coordinatePolarization_one_set_neg_one_imp}
\label{FABL.coalitionCanAttain_coordinatePolarization_one_set_neg_one_imp}
\leanok
\uses{FABL.BooleanFunction, FABL.CoalitionCanAttain, FABL.Sign, FABL.SignCube, FABL.coordinatePolarization, FABL.setCoordinate, FABL.setCoordinate_apply_of_ne, FABL.sign_le_one}
\difficulty{5}
Let $n$ be a natural number, let $f\colon\{−1,1\}^n\to\{−1,1\}$, let
$J\subseteq\{1,\ldots,n\}$, and choose a coordinate $i\notin J$ and an input
$x\in\{−1,1\}^n$. Define the coordinate-polarized function $f_i^{\mathrm{pol}}$ by
taking, at an input $z$, the maximum of $f$ at the two inputs obtained by setting
coordinate $i$ to $1$ and $-1$ if $z_i=1$, and their minimum if $z_i=-1$. If
$f_i^{\mathrm{pol}}$ can attain $1$ at an input agreeing with $x$ with coordinate $i$
set to $-1$, and agreeing with that input outside $J$, then $f$ can attain $1$ from an
input with coordinate $i$ set to $1$ or from one with coordinate $i$ set to $-1$,
allowing changes only in $J$ from the respective input.
\end{theorem}

\begin{theorem}[Attainability from both settings of a coordinate]
\lean{FABL.coalitionCanAttain_coordinatePolarization_neg_one_set_one_imp}
\label{FABL.coalitionCanAttain_coordinatePolarization_neg_one_set_one_imp}
\leanok
\uses{FABL.BooleanFunction, FABL.CoalitionCanAttain, FABL.Sign, FABL.SignCube, FABL.coordinatePolarization, FABL.neg_one_le_sign, FABL.setCoordinate, FABL.setCoordinate_apply_of_ne}
\difficulty{5}
Let $n$ be a natural number, let $f\colon\{−1,1\}^n\to\{−1,1\}$, let
$J\subseteq\{1,\ldots,n\}$, and choose a coordinate $i\notin J$ and an input
$x\in\{−1,1\}^n$. Define the coordinate-polarized function $f_i^{\mathrm{pol}}$ by
taking, at an input $z$, the maximum of $f$ at the two inputs obtained by setting
coordinate $i$ to $1$ and $-1$ if $z_i=1$, and their minimum if $z_i=-1$. If
$f_i^{\mathrm{pol}}$ can attain $-1$ at an input agreeing with $x$ with coordinate $i$
set to $1$, and agreeing with that input outside $J$, then $f$ can attain $-1$ from
inputs with coordinate $i$ set to $1$ and with coordinate $i$ set to $-1$, in both cases
allowing changes only in $J$ from the corresponding input.
\end{theorem}

\begin{theorem}[Attainability after setting a coordinate to minus one]
\lean{FABL.coalitionCanAttain_coordinatePolarization_neg_one_set_neg_one_imp}
\label{FABL.coalitionCanAttain_coordinatePolarization_neg_one_set_neg_one_imp}
\leanok
\uses{FABL.BooleanFunction, FABL.CoalitionCanAttain, FABL.Sign, FABL.SignCube, FABL.coordinatePolarization, FABL.setCoordinate, FABL.setCoordinate_apply_of_ne}
\difficulty{5}
Let $n$ be a natural number, let $f\colon\{−1,1\}^n\to\{−1,1\}$, let
$J\subseteq\{1,\ldots,n\}$, and choose a coordinate $i\notin J$ and an input
$x\in\{−1,1\}^n$. Define the coordinate-polarized function $f_i^{\mathrm{pol}}$ by
taking, at an input $z$, the maximum of $f$ at the two inputs obtained by setting
coordinate $i$ to $1$ and $-1$ if $z_i=1$, and their minimum if $z_i=-1$. If
$f_i^{\mathrm{pol}}$ can attain $-1$ at an input agreeing with $x$ with coordinate $i$
set to $-1$, and agreeing with that input outside $J$, then $f$ can attain $-1$ at an
input agreeing with $x$ with coordinate $i$ set to $1$, or at one agreeing with $x$ with
coordinate $i$ set to $-1$; in each case agreement is required outside $J$.
\end{theorem}

\begin{theorem}[Monotonicity of the mean]
\lean{FABL.mean_le_mean_of_pointwise}
\label{FABL.mean_le_mean_of_pointwise}
\leanok
\uses{FABL.SignCube, FABL.mean}
\difficulty{3}
Let $n$ be a nonnegative integer, and let $u,v:\{−1,1\}^n\to\mathbb R$ satisfy $u(x)\le
v(x)$ for every input $x$. Then the uniform mean of $u$ is at most the uniform mean of
$v$.
\end{theorem}

\begin{theorem}[Indicator bound from union and intersection]
\lean{FABL.indicator_pair_le_of_union_inter}
\label{FABL.indicator_pair_le_of_union_inter}
\leanok
\difficulty{4}
Let $P_+,P_-,A_+,A_-$ be propositions such that $P_+$ implies $A_+\vee A_-$ and $P_-$
implies $A_+\wedge A_-$. Then $\mathbf{1}_{P_+}+\mathbf{1}_{P_-}\leq
\mathbf{1}_{A_+}+\mathbf{1}_{A_-}$, where each indicator is $1$ when its proposition
holds and $0$ otherwise.
\end{theorem}

\begin{theorem}[Monotonicity of a pair of indicators]
\lean{FABL.indicator_pair_le_of_imp}
\label{FABL.indicator_pair_le_of_imp}
\leanok
\difficulty{4}
Let $P_+,P_-,A_+,A_-$ be decidable propositions. If $P_+$ implies $A_+$ and $P_-$
implies $A_-$, then $\mathbf 1_{P_+}+\mathbf 1_{P_-}\leq\mathbf 1_{A_+}+\mathbf
1_{A_-}$, where each indicator is $1$ when its proposition holds and $0$ otherwise.
\end{theorem}

\begin{theorem}[Coalition Success Under Coordinate Polarization]
\lean{FABL.coalitionSuccessProbability_coordinatePolarization_le}
\label{FABL.coalitionSuccessProbability_coordinatePolarization_le}
\leanok
\uses{FABL.BooleanFunction, FABL.CoalitionCanAttain, FABL.Sign, FABL.coalitionCanAttain_coordinatePolarization_imp_of_mem, FABL.coalitionCanAttain_coordinatePolarization_neg_one_set_neg_one_imp, FABL.coalitionCanAttain_coordinatePolarization_neg_one_set_one_imp, FABL.coalitionCanAttain_coordinatePolarization_one_set_neg_one_imp, FABL.coalitionCanAttain_coordinatePolarization_one_set_one_imp, FABL.coalitionSuccessProbability, FABL.coordinateExpectation_apply, FABL.coordinatePolarization, FABL.indicator_pair_le_of_imp, FABL.indicator_pair_le_of_union_inter, FABL.mean, FABL.mean_coordinateExpectation, FABL.mean_le_mean_of_pointwise, FABL.setCoordinate, FABL.uniformProbability}
\difficulty{6}
Let $n$ be a nonnegative integer, let $f:\{-1,1\}^n\to\{-1,1\}$, let $J$ be a set of
coordinates, let $b\in\{-1,1\}$, and choose a coordinate $i$. Form the coordinate
polarization of $f$ by taking, at inputs with $i$th coordinate $1$, the maximum of the
values of $f$ at the two inputs with coordinate $i$ set to $1$ and $-1$, respectively,
and taking their minimum when the $i$th coordinate is $-1$. The fraction of inputs from
which the coalition $J$ can make the function attain $b$ for this polarized function is
at most that fraction for $f$.
\end{theorem}

\begin{theorem}[Coalitional Influence Under Coordinate Polarization]
\lean{FABL.coalitionalInfluenceToward_coordinatePolarization_le}
\label{FABL.coalitionalInfluenceToward_coordinatePolarization_le}
\leanok
\uses{FABL.BooleanFunction, FABL.Sign, FABL.coalitionSuccessProbability_coordinatePolarization_le, FABL.coalitionalInfluenceToward, FABL.coordinatePolarization, FABL.uniformProbability_coordinatePolarization_eq}
\difficulty{3}
Let $n$ be a nonnegative integer, let $f:\{-1,1\}^n\to\{-1,1\}$, let $J$ be a set of
coordinates, let $b\in\{-1,1\}$, and choose a coordinate $i$. Form the coordinate
polarization of $f$ by taking, at inputs with $i$th coordinate $1$, the maximum of the
values of $f$ at the two inputs with coordinate $i$ set to $1$ and $-1$, respectively,
and taking their minimum when the $i$th coordinate is $-1$. Define the influence toward
$b$ as the fraction of inputs from which $J$ can make the function attain $b$ minus the
fraction of inputs where the function already equals $b$. The influence toward $b$ of
the polarized function is at most that of $f$.
\end{theorem}

\begin{theorem}[Polarization decreases coalitional influence]
\lean{FABL.coalitionalInfluence_coordinatePolarization_le}
\label{FABL.coalitionalInfluence_coordinatePolarization_le}
\leanok
\uses{FABL.BooleanFunction, FABL.coalitionalInfluence, FABL.coalitionalInfluenceToward_coordinatePolarization_le, FABL.coalitionalInfluence_eq_toward_add_toward, FABL.coordinatePolarization}
\difficulty{3}
Let $f:\{−1,1\}^n\to\{−1,1\}$ be a Boolean function, let $J$ be a set of coordinates,
and let $i$ be a coordinate. Form the polarization of $f$ at $i$ by assigning, at each
input, the maximum of the two values obtained by setting coordinate $i$ to $1$ or $-1$
when that input has coordinate $i=1$, and the minimum of those values when it has
coordinate $i=-1$. The coalitional influence of this polarized function on $J$ is at
most that of $f$ on $J$.
\end{theorem}

\begin{definition}[Coordinate polarization]
\lean{FABL.polarizeCoordinates}
\statementsource{Analysis-of-Boolean-Functions-by-Ryan-ODonnell}{pdf-ffbe9c7b3306-p065-b005}
\label{FABL.polarizeCoordinates}
\leanok
\uses{FABL.BooleanFunction, FABL.coordinatePolarization}
Let $f:\{-1,1\}^n\to\{\pm1\}$ and let $i$ be a coordinate. The polarization of $f$ at
$i$ takes, at an input with $i$th coordinate $1$, the maximum of the two values of $f$
obtained by setting that coordinate to $1$ and $-1$; at an input with $i$th coordinate
$-1$, it takes their minimum. For a finite list of coordinates, successive polarization
means applying these operations in list order, with the empty list leaving $f$
unchanged.
\end{definition}

\begin{theorem}[Polarization along the empty list]
\lean{FABL.polarizeCoordinates_nil}
\label{FABL.polarizeCoordinates_nil}
\leanok
\uses{FABL.BooleanFunction, FABL.polarizeCoordinates}
\difficulty{1}
For every natural number $n$ and every function from the $n$-dimensional sign cube to
$\{ -1,1\}$, polarizing successively along the empty list of coordinates returns the
original function.
\end{theorem}

\begin{theorem}[Polarization along a list]
\lean{FABL.polarizeCoordinates_cons}
\label{FABL.polarizeCoordinates_cons}
\leanok
\uses{FABL.BooleanFunction, FABL.coordinatePolarization, FABL.polarizeCoordinates}
\difficulty{1}
For every Boolean function on the $n$-dimensional sign cube, coordinate $i$, and finite
list of coordinates $L$, polarizing along the list beginning with $i$ is the same as
first polarizing the function at coordinate $i$ and then polarizing along $L$.
\end{theorem}

\begin{theorem}[Mean preserved by coordinate polarization]
\lean{FABL.mean_polarizeCoordinates}
\proofsource{Analysis-of-Boolean-Functions-by-Ryan-ODonnell}{pdf-ffbe9c7b3306-p065-b005}
\label{FABL.mean_polarizeCoordinates}
\leanok
\uses{FABL.BooleanFunction, FABL.BooleanFunction.toReal, FABL.mean, FABL.mean_coordinatePolarization, FABL.polarizeCoordinates, FABL.polarizeCoordinates_cons}
\difficulty{4}
Let $f$ be a function from the $n$-dimensional sign cube to $\{-1,1\}$, and let $I$ be any
finite list of coordinates, with repetitions allowed. After successively polarizing $f$
in the coordinates listed in $I$, the mean of the resulting real-valued function equals
the mean of $f$ regarded as real-valued.
\end{theorem}

\begin{theorem}[Influence Under Coordinate Polarization]
\lean{FABL.coalitionalInfluenceToward_polarizeCoordinates_le}
\label{FABL.coalitionalInfluenceToward_polarizeCoordinates_le}
\leanok
\uses{FABL.BooleanFunction, FABL.Sign, FABL.coalitionalInfluenceToward, FABL.coalitionalInfluenceToward_coordinatePolarization_le, FABL.coordinatePolarization, FABL.polarizeCoordinates, FABL.polarizeCoordinates_cons}
\difficulty{4}
Let $n$ be a nonnegative integer, let $f:\{-1,1\}^n\to\{-1,1\}$, and let
$(i_1,\ldots,i_k)$ be any finite list of coordinates in $\{0,\ldots,n-1\}$. Obtain
$f^{(i_1,\ldots,i_k)}$ by successively polarizing in those coordinates: at coordinate
$i$, the new value at $x$ is the larger of the two values obtained by setting $x_i$ to
$+1$ or $-1$ when $x_i=+1$, and the smaller when $x_i=-1$. For every coordinate set
$J\subseteq\{0,\ldots,n-1\}$ and target sign $b$, the coalitional influence of
$f^{(i_1,\ldots,i_k)}$ toward $b$ is at most that of $f$ toward $b$.
\end{theorem}

\begin{theorem}[Polarization reduces coalitional influence]
\lean{FABL.coalitionalInfluence_polarizeCoordinates_le}
\label{FABL.coalitionalInfluence_polarizeCoordinates_le}
\leanok
\uses{FABL.BooleanFunction, FABL.coalitionalInfluence, FABL.coalitionalInfluence_coordinatePolarization_le, FABL.coordinatePolarization, FABL.polarizeCoordinates, FABL.polarizeCoordinates_cons}
\difficulty{4}
Let $n$ be a nonnegative integer, let $f$ be a function from the sign cube $\{-1,1\}^n$
to $\{-1,1\}$, let $(i_1,\ldots,i_m)$ be a finite list of coordinates, and let $J$ be a
set of coordinates. The coalitional influence of the function obtained by successively
polarizing $f$ in the listed coordinates is at most the coalitional influence of $f$ on
$J$.
\end{theorem}

\begin{theorem}[Polarization sequences preserve monotonicity]
\lean{FABL.isMonotoneInCoordinate_polarizeCoordinates}
\proofsource{Analysis-of-Boolean-Functions-by-Ryan-ODonnell}{pdf-ffbe9c7b3306-p065-b005}
\label{FABL.isMonotoneInCoordinate_polarizeCoordinates}
\leanok
\uses{FABL.BooleanFunction, FABL.IsMonotoneInCoordinate, FABL.coordinatePolarization, FABL.isMonotoneInCoordinate_coordinatePolarization, FABL.polarizeCoordinates, FABL.polarizeCoordinates_cons}
\difficulty{4}
Let $n$ be a nonnegative integer, let $f$ be a sign-valued function on the
$n$-dimensional Boolean cube, let $is$ be any finite list of coordinates, and let $j$ be
a coordinate. If $f$ is monotone in coordinate $j$, then the function obtained by
successively polarizing $f$ along the list $is$ remains monotone in coordinate $j$.
\end{theorem}

\begin{theorem}[Listed coordinate becomes monotone]
\lean{FABL.isMonotoneInCoordinate_polarizeCoordinates_of_mem}
\label{FABL.isMonotoneInCoordinate_polarizeCoordinates_of_mem}
\leanok
\uses{FABL.BooleanFunction, FABL.IsMonotoneInCoordinate, FABL.coordinatePolarization, FABL.isMonotoneInCoordinate_coordinatePolarization_self, FABL.isMonotoneInCoordinate_polarizeCoordinates, FABL.polarizeCoordinates, FABL.polarizeCoordinates_cons}
\difficulty{5}
Let $n$ be a nonnegative integer, let $f$ be a sign-valued function on the
$n$-dimensional Boolean cube, let $is$ be any finite list of coordinates, and let $j$ be
a coordinate appearing in $is$. After successively polarizing $f$ along $is$, the
resulting function is monotone in coordinate $j$.
\end{theorem}

\begin{theorem}[Monotonicity under coordinate bribery]
\lean{FABL.le_bribeInput_one_of_isMonotoneInCoordinate}
\label{FABL.le_bribeInput_one_of_isMonotoneInCoordinate}
\leanok
\uses{FABL.BooleanFunction, FABL.IsMonotoneInCoordinate, FABL.SignCube, FABL.bribeInput, FABL.bribeInput_insert, FABL.setCoordinate, FABL.setCoordinate_eq_self}
\difficulty{5}
Let $f:\{−1,1\}^n\to\{−1,1\}$ be nondecreasing in each coordinate, meaning that changing
any one coordinate from $-1$ to $1$ cannot decrease its value. For every finite set of
coordinates $J\subseteq\{1,\ldots,n\}$ and every input $x$, the value of $f$ at $x$ is
at most its value at the input obtained by setting every coordinate in $J$ to $1$ and
leaving all other coordinates unchanged.
\end{theorem}

\begin{theorem}[Coordinatewise monotonicity implies monotonicity]
\lean{FABL.monotone_of_isMonotoneInCoordinate}
\label{FABL.monotone_of_isMonotoneInCoordinate}
\leanok
\uses{FABL.BooleanFunction, FABL.IsMonotoneInCoordinate, FABL.Sign, FABL.bribeInput, FABL.le_bribeInput_one_of_isMonotoneInCoordinate, FABL.neg_one_le_sign}
\difficulty{5}
Let $f:\{-1,1\}^n\to\{-1,1\}$. If, for every coordinate $i$ and every input $x$,
changing coordinate $i$ from $-1$ to $1$ does not decrease the value of $f$, then $f$ is
monotone on the sign cube.
\end{theorem}

\begin{definition}[Coordinate monotonization]
\lean{FABL.monotonization}
\statementsource{Analysis-of-Boolean-Functions-by-Ryan-ODonnell}{
  pdf-ffbe9c7b3306-p065-b005,
  pdf-ffbe9c7b3306-p275-b006
}
\label{FABL.monotonization}
\leanok
\uses{FABL.BooleanFunction, FABL.polarizeCoordinates}
For a Boolean function $f:\{-1,1\}^n\to\{\pm1\}$, its monotonization is obtained by
successively polarizing in every coordinate, in the order $1,2,\ldots,n$.
\end{definition}

\begin{theorem}[Mean preserved by monotonization]
\lean{FABL.mean_monotonization}
\label{FABL.mean_monotonization}
\leanok
\uses{FABL.BooleanFunction, FABL.BooleanFunction.toReal, FABL.mean, FABL.mean_polarizeCoordinates, FABL.monotonization}
\difficulty{1}
Let $n$ be a nonnegative integer and let $f:\{−1,1\}^n\to\{−1,1\}$. The uniform mean of
the real-valued function obtained by monotonizing $f$ equals the uniform mean of $f$,
with its values viewed as real numbers.
\end{theorem}

\begin{theorem}[Monotonization is coordinatewise monotone]
\lean{FABL.isMonotoneInCoordinate_monotonization}
\proofsource{Analysis-of-Boolean-Functions-by-Ryan-ODonnell}{pdf-ffbe9c7b3306-p065-b005}
\label{FABL.isMonotoneInCoordinate_monotonization}
\leanok
\uses{FABL.BooleanFunction, FABL.IsMonotoneInCoordinate, FABL.isMonotoneInCoordinate_polarizeCoordinates_of_mem, FABL.monotonization}
\difficulty{3}
Let $n$ be a nonnegative integer and let $f$ be a sign-valued function on the
$n$-dimensional Boolean cube. The function obtained by successively polarizing $f$ in
every coordinate is monotone in each coordinate.
\end{theorem}

\begin{theorem}[Monotonicity of coordinate polarization]
\lean{FABL.monotone_monotonization}
\proofsource{Analysis-of-Boolean-Functions-by-Ryan-ODonnell}{pdf-ffbe9c7b3306-p065-b005}
\label{FABL.monotone_monotonization}
\leanok
\uses{FABL.BooleanFunction, FABL.isMonotoneInCoordinate_monotonization, FABL.monotone_of_isMonotoneInCoordinate, FABL.monotonization}
\difficulty{1}
For every function $f:\{ -1,1\}^n\to\{-1,1\}$, polarizing its coordinates successively
in the order $1,\ldots,n$ produces a monotone function on the sign cube.
\end{theorem}

\begin{theorem}[Coalitional influence under monotonization]
\lean{FABL.coalitionalInfluenceToward_monotonization_le}
\proofsource{Analysis-of-Boolean-Functions-by-Ryan-ODonnell}{pdf-ffbe9c7b3306-p275-b006}
\label{FABL.coalitionalInfluenceToward_monotonization_le}
\leanok
\uses{FABL.BooleanFunction, FABL.Sign, FABL.coalitionalInfluenceToward, FABL.coalitionalInfluenceToward_polarizeCoordinates_le, FABL.monotonization}
\difficulty{1}
Let $n$ be a natural number, let $f$ be a function from the $n$-dimensional sign cube to
the signs, let $J$ be a set of coordinates, and let $b$ be a sign. The coalitional
influence of the coordinatewise monotonization of $f$ toward $b$ is at most the
coalitional influence of $f$ toward $b$.
\end{theorem}

\begin{theorem}[Monotonization reduces coalitional influence]
\lean{FABL.coalitionalInfluence_monotonization_le}
\proofsource{Analysis-of-Boolean-Functions-by-Ryan-ODonnell}{pdf-ffbe9c7b3306-p275-b006}
\label{FABL.coalitionalInfluence_monotonization_le}
\leanok
\uses{FABL.BooleanFunction, FABL.coalitionalInfluence, FABL.coalitionalInfluence_polarizeCoordinates_le, FABL.monotonization}
\difficulty{1}
Let $n$ be a nonnegative integer, let $f$ be a function from the sign cube
$\{-1,1\}^n$ to $\{-1,1\}$, and let $J$ be a set of coordinates. The coalitional
influence of the coordinatewise monotonization of $f$ on $J$ is at most the coalitional
influence of $f$ on $J$.
\end{theorem}

\begin{theorem}[Variance under coordinate monotonization]
\lean{FABL.variance_monotonization}
\label{FABL.variance_monotonization}
\leanok
\uses{FABL.BooleanFunction, FABL.BooleanFunction.toReal, FABL.mean_monotonization, FABL.monotonization, FABL.variance, FABL.variance_eq_four_mul_probabilities}
\difficulty{3}
Let $n$ be a natural number and let $f:\{−1,1\}^n\to\{−1,1\}$ be a Boolean function. The
variance of the real-valued function obtained by monotonizing $f$ coordinate by
coordinate equals the variance of $f$ viewed as a real-valued function.
\end{theorem}

\begin{theorem}[Coalition influence bound]
\lean{FABL.exercise9_28_coalition}
\proofsource{Analysis-of-Boolean-Functions-by-Ryan-ODonnell}{
  pdf-ffbe9c7b3306-p275-b005,
  pdf-ffbe9c7b3306-p276-b001
}
\label{FABL.exercise9_28_coalition}
\leanok
\uses{FABL.BooleanFunction, FABL.BooleanFunction.toReal, FABL.coalitionalInfluence, FABL.coalitionalInfluence_monotonization_le, FABL.exercise9_27_coalition, FABL.monotone_monotonization, FABL.monotonization, FABL.variance, FABL.variance_monotonization}
\difficulty{4}
Let $f:\{-1,1\}^n\to\{-1,1\}$, with $2\le n$. Suppose $\delta>0$ and $\delta$ is at most
the variance of the real-valued sign representation of $f$. For $0<\varepsilon<1/2$,
there is a coordinate set $J$ whose coalitional influence is at least $1-\varepsilon$
and whose size satisfies $|J|\le 2+\frac{200n}{\log
n}\log\!\left(\frac{8}{\varepsilon\delta}\right)$.
\end{theorem}
